\section{With Their Own Bodies}

Hell is not only the state of a separated soul. The Athanasian Creed says that at Christ's coming ``all men shall rise again with their bodies,''\endnote{\latin{Quicumque}, DS 76, in the English of the \emph{Catholic Encyclopedia} (1907).} and the Fourth Lateran Council professed that all will rise with their own bodies which they now bear, \latin{cum suis propriis resurgent corporibus, quae nunc gestant}, to receive according to their works.\endnote{Fourth Lateran Council, \latin{Firmiter}, DS 801.} The resurrection is a gift of Christ's victory, and it is given to all; but for those who have refused him it becomes the occasion of a fuller loss. The body that was the soul's partner in sin rises to be its partner in punishment.

The Fathers taught this from the beginning. Saint Justin Martyr told the Roman emperor that Christ will judge the wicked ``in the same bodies united again to their spirits which are now to undergo everlasting punishment; and not only, as Plato said, for a period of a thousand years.''\endnote{Justin Martyr, \emph{First Apology} 8, \emph{The Ante-Nicene Fathers}, vol.~1.} Cyprian wrote: ``Souls with their bodies will be reserved in infinite tortures for suffering.''\endnote{Cyprian, \emph{To Demetrianus} 24.} And Saint Cyril of Jerusalem, preparing catechumens for baptism, set the two resurrections side by side and gave the reason in a sentence a child could remember:

\begin{quote}
We shall be raised therefore, all with our bodies eternal, but not all with bodies alike: for if a man is righteous, he will receive a heavenly body, that he may be able worthily to hold converse with Angels; but if a man is a sinner, he shall receive an eternal body, fitted to endure the penalties of sins, that he may burn eternally in fire, nor ever be consumed. And righteously will God assign this portion to either company; for we do nothing without the body. We blaspheme with the mouth, and with the mouth we pray. With the body we commit fornication, and with the body we keep chastity. With the hand we rob, and by the hand we bestow alms; and the rest in like manner. Since then the body has been our minister in all things, it shall also share with us in the future the fruits of the past.\endnote{Cyril of Jerusalem, \emph{Catechetical Lectures} 18.19, trans.\ E. H. Gifford, \emph{Nicene and Post-Nicene Fathers}, 2nd series, vol.~7.}
\end{quote}

\noindent This is why Augustine and the Supplement, whatever they allowed about the worm, held the fire that torments the risen flesh to be bodily.

The doctrine is a consolation before it is a warning. It means that the body matters: that the hand, the eye, the tongue, the flesh are not instruments a man uses and lays aside, but parts of himself that will stand with him before God. The same mouth that prays will rise; the same hands that give alms will rise. Cyril draws the lesson at once: ``Therefore, brethren, let us be careful of our bodies, nor misuse them as though not our own.''\endnote{Cyril, \emph{Catechetical Lectures} 18.20.}

It also explains the strange severity of Christ's counsel to cut off the hand and pluck out the eye. Augustine reads the members as ``those persons whom a man loves as the most useful members of his body'';\endnote{Augustine, \emph{City of God} XXI.9.} the Lord commands a man to treat whatever leads him into sin as he would treat a limb that threatened his life. The occasion that seems as dear as an eye---the friendship, the habit, the pleasure, the place---is to be cut away, because the whole body is at stake, and the whole soul with it. ``And fear ye not them that kill the body, and are not able to kill the soul: but rather fear him that can destroy both soul and body in hell.''\endnote{Matthew 10:28, Douay--Rheims (Challoner).} The body that is kept from sin now, by fasting and watchfulness and the avoidance of occasions, is being kept for the resurrection of life.
